\documentclass[a4paper,11pt]{article}
\usepackage[utf8]{inputenc}
\usepackage{geometry}
\usepackage{enumitem}
\usepackage{titlesec}
\usepackage{xcolor}
\usepackage{hyperref}
\usepackage{fancyhdr}

% Page Setup
\geometry{top=1in, bottom=1in, left=0.8in, right=0.8in}
\pagestyle{fancy}
\fancyhf{}
\lhead{\small Project 83: Agentic AI Compiler Assistant}
\rhead{\small Execution Plan 2026}
\cfoot{\thepage}

% Formatting
\titleformat{\section}{\large\bfseries\color{blue!70!black}}{Phase \thesection:}{1em}{}[\titlerule]
\titleformat{\subsection}{\bfseries}{Week \thesubsection:}{1em}{}
\newcommand{\weekdetail}[3]{
    \textbf{Tasks:} #1 \\
    \textbf{Tools:} #2 \\
    \textbf{Resources/Methods:} #3 \\
}

\begin{document}

\begin{center}
    {\huge \textbf{Detailed Execution Plan and Report}} \\
    \vspace{2mm}
    {\large \textbf{Project 83: Conversational Agentic AI Compiler Assistant}} \\
    \vspace{2mm}
    \textit{Target Language: Python (Subset) | Intelligence: Google Gemini Pro API}
\end{center}

\section{Problem Understanding \& Basic Foundations}

\subsection{1 -- Problem Definition \& Scope Research}
\weekdetail{
    Define the EBNF (Extended Backus-Naur Form) grammar for a Python subset (Mini-Python) including variables, loops, and functions. Map the "Agentic" interaction flow.
}{
    LucidChart (Grammar Mapping), Markdown (Documentation).
}{
    Python Language Reference (Grammar Specification), Research on "ReAct" (Reason + Act) agent patterns.
}

\subsection{2 -- Gemini API \& Agentic Workflow Design}
\weekdetail{
    Set up Google AI Studio environment. Test Gemini's capability to interpret raw Python tokens and AST structures. Design the prompt architecture for the "Auditor" and "Assistant" agents.
}{
    Google Generative AI SDK, Python `python-dotenv`.
}{
    Gemini API Documentation, Prompt Engineering Guide for Developers (DeepLearning.AI).
}

\section{Requirement Analysis \& System Design}

\subsection{3 -- Requirements \& Security Policy Analysis}
\weekdetail{
    Identify functional requirements (error detection, code fix suggestions) and security policies (detecting `eval()`, insecure imports, or infinite recursion).
}{
    Jira/Trello (Task Tracking), LaTeX (SRS Document).
}{
    OWASP Top 10 for LLM Applications, Python Security Best Practices.
}

\subsection{4 -- System Architecture \& Orchestration Planning}
\weekdetail{
    Design the "Compiler-Agent Bridge." Define how the compiler state (AST) is converted to JSON for Gemini and how Gemini's responses trigger compiler re-runs.
}{
    Draw.io (Architecture Diagrams), Mermaid.js.
}{
    "Building LLM-Powered Agents" (NVIDIA Technical Blog).
}

\section{Compiler Front-End Implementation (From Scratch)}

\subsection{5 -- Lexer (Lexical Analyzer) Development}
\weekdetail{
    Build a Regex-based Lexer from scratch. Convert source code into a stream of tokens (Keywords, Identifiers, Literals, Operators). Handle indentation tracking (unique to Python).
}{
    Python `re` module, `enum` for Token Types.
}{
    "Compilers: Principles, Techniques, and Tools" (Dragon Book - Chapter 3).
}

\subsection{6 -- Syntax Parser (Recursive Descent)}
\weekdetail{
    Implement a Top-Down Recursive Descent Parser. Construct an Abstract Syntax Tree (AST) from the token stream. Handle operator precedence and nested blocks.
}{
    Python classes (Node objects), `json` (for tree serialization).
}{
    "Crafting Interpreters" by Robert Nystrom (Parsing section).
}

\subsection{7 -- Semantic Analyzer \& Basic Diagnostics}
\weekdetail{
    Implement a Symbol Table to track variable scope and type checking. Build a diagnostic engine that logs errors with line and column numbers.
}{
    Python Dictionary-based Symbol Tables.
}{
    Modern Compiler Implementation in ML (Appel).
}

\section{AI Integration \& Agentic Logic}

\subsection{8 -- Gemini API Integration \& Context Injection}
\weekdetail{
    Develop the "Context Provider" that formats the AST and Symbol Table into a natural language description for Gemini. Implement secure API call handling.
}{
    Gemini-1.5-Pro API, `google-generativeai` library.
}{
    JSON-to-Text serialization techniques for LLMs.
}

\subsection{9 -- Agentic Tool-Use (Function Calling)}
\weekdetail{
    Enable Gemini to "call" compiler functions (e.g., `check_scope()`, `reparse_code()`) to investigate errors autonomously before responding to the user.
}{
    Gemini Function Calling API.
}{
    "Model-controlled tool use" (Google AI Documentation).
}

\subsection{10 -- Conversational Debugging Interface}
\weekdetail{
    Develop the Chatbot UI/CLI where the user submits code. The AI should explain errors "step-by-step" rather than just outputting a traceback.
}{
    Streamlit (Web UI) or Rich (CLI).
}{
    Human-Computer Interaction (HCI) principles for AI assistants.
}

\section{Security, Optimization \& Testing}

\subsection{11 -- Security Risk Auditing Agent}
\weekdetail{
    Program a specific Agent role to audit code for "Logic Bombs" or "Prompt Injection" within the code comments. Ensure the compiler rejects unsafe constructs.
}{
    Static Analysis Security Testing (SAST) logic.
}{
    Snyk Python Security scan reports.
}

\subsection{12 -- Performance Benchmarking}
\weekdetail{
    Measure the latency of AI responses. Compare the accuracy of Gemini's suggested fixes against traditional Python `SyntaxError` messages.
}{
    Python `timeit` module, `PyTest` for automated testing.
}{
    Standard Python error test suites.
}

\section{Refinement \& Presentation}

\subsection{13 -- Explainability \& Traceability (XAI)}
\weekdetail{
    Implement a "Reasoning Log" where the agent shows which part of the AST led to its decision. This fulfills the Explainability requirement of the project.
}{
    Graphviz (AST visualization).
}{
    Explainable AI (XAI) framework standards.
}

\subsection{14 -- Integration & Bug Squashing}
\weekdetail{
    Perform end-to-end testing. Ensure that code fixed by the AI actually passes the "from scratch" parser without errors.
}{
    GitHub Actions (CI/CD).
}{
    Regression testing methodologies.
}

\subsection{15 -- Documentation \& Report Finalization}
\weekdetail{
    Complete the final report including architecture diagrams, prompt templates, and code snippets. Prepare presentation slides.
}{
    Overleaf (LaTeX), Google Slides.
}{
    Academic Report Writing standards.
}

\subsection{16 -- Live Demonstration \& Submission}
\weekdetail{
    Conduct a live demo of the Conversational Agentic Assistant. Submit the final source code repository and documentation.
}{
    Loom (Demo Recording), GitHub (Code Repository).
}{
    Project Rubric and Assessment Criteria.
}

\end{document}